\documentclass[11pt,a4paper]{article}

\usepackage[spanish]{babel}
\usepackage[T1]{fontenc}
\usepackage[utf8]{inputenc}
\usepackage{lmodern}
\usepackage{geometry}
\usepackage{setspace}
\usepackage{booktabs}
\usepackage{longtable}
\usepackage{array}
\usepackage{hyperref}
\usepackage{graphicx}
\usepackage{xcolor}
\usepackage{amsmath}
\usepackage{siunitx}
\usepackage{listings}
\usepackage{enumitem}

\geometry{margin=2.5cm}
\setstretch{1.12}
\hypersetup{
  colorlinks=true,
  linkcolor=blue,
  urlcolor=blue,
  citecolor=blue,
  pdftitle={Memoria del proyecto Rescue Video Analyzer},
  pdfauthor={Proyecto de Arquitectura Paralela}
}

\lstdefinelanguage{bash}{
  morekeywords={cmake,mpirun,python,wget,curl,sudo,apt,get,install,export},
  sensitive=true,
  morecomment=[l]{\#},
  morestring=[b]"
}

\title{Memoria t\'ecnica completa\\Rescue Video Analyzer}
\author{Arquitectura de Computadores y Programaci\'on Paralela}
\date{\today}

\begin{document}

\maketitle
\tableofcontents
\newpage

\section{Resumen ejecutivo}
Este proyecto implementa un analizador de video de personas sobre una arquitectura distribuida y \textbf{altamente heterog\'enea} que combina \textbf{MPI}, \textbf{CUDA (NVIDIA)} y \textbf{OpenCL (Intel/AMD)}. El sistema usa un esquema maestro-trabajadores donde el \emph{rank} 0 lee el video, serializa lotes de frames y los env\'ia por MPI; los \emph{workers} reciben esos frames, los preprocesan nativamente en el hardware gr\'afico disponible y ejecutan inferencia con YOLO11 aprovechando la m\'axima capacidad de c\'omputo del sistema, combinando GPUs dedicadas e integradas de forma simult\'anea.

Los resultados se consolidan en tres salidas:
\begin{itemize}[leftmargin=*]
  \item video anotado con cajas y overlay informativo,
  \item CSV por frame con m\'etricas y conteos,
  \item JSON resumen con estad\'isticas globales y balance de carga.
\end{itemize}

Como mejora funcional reciente, el video anotado incluye:
\begin{itemize}[leftmargin=*]
  \item \textbf{Personas actuales}: personas detectadas en el frame actual.
  \item \textbf{Acumulado real}: estimaci\'on de personas \`unicas a partir de IDs de tracking visual.
\end{itemize}

\section{Objetivos del proyecto}
\begin{enumerate}[leftmargin=*]
  \item Dise\~nar un pipeline de procesamiento distribuido de video con MPI.
  \item Acelerar preprocesado e inferencia con GPU CUDA.
  \item Generar m\'etricas reproducibles de rendimiento por etapa.
  \item Mantener trazabilidad de resultados a nivel de frame y a nivel global.
  \item Producir evidencia visual del an\'alisis mediante video anotado.
\end{enumerate}

\section{Alcance y restricciones}
\subsection{Alcance}
\begin{itemize}[leftmargin=*]
  \item Detecci\'on de personas en video con modelos YOLO11 ONNX.
  \item Orquestaci\'on de trabajo con MPI permitiendo ejecuci\'on multiproceso.
  \item Ejecuci\'on polim\'orfica escalable horizontalmente, utilizando m\'ultiples tarjetas gr\'aficas f\'isicas de distintos fabricantes (NVIDIA, Intel, AMD) en paralelo.
\end{itemize}

\subsection{Restricciones principales}
\begin{itemize}[leftmargin=*]
  \item El tracker se usa para estabilizar la visualizaci\'on, no para garantizar una identidad robusta a trav\'es de oclusiones prolongadas.
  \item El conteo acumulado es una estimaci\'on basada en IDs de tracking visual y puede variar en escenarios de aglomeraci\'on densa.
  \item Requiere la instalaci\'on de controladores OpenCL (ICD) en el sistema operativo host para reconocer las GPUs integradas de Intel/AMD.
\end{itemize}

\subsection{Entorno exacto (probado)}
\begin{itemize}[leftmargin=*]
  \item El binario busca bibliotecas CUDA/cuDNN dentro de una \path{.venv} en rutas de Python 3.12 o 3.11, y pre-carga librer\'ias cu13/cuDNN antes de crear la sesi\'on ORT.
  \item Si existe \path{/usr/lib/wsl/lib}, se agrega autom\'aticamente al \texttt{LD\_LIBRARY\_PATH} para soporte WSL.
  \item ONNX Runtime CUDA se carga din\'amicamente desde la \path{.venv}; si es necesario se puede forzar la librer\'ia con \texttt{RESCUE\_ORT\_LIBRARY}.
  \item La compilaci\'on CUDA usa \texttt{CMAKE\_CUDA\_ARCHITECTURES=89} por defecto (se puede sobrescribir al invocar CMake).
\end{itemize}

\section{Arquitectura del sistema}
\subsection{Vista general}
\begin{enumerate}[leftmargin=*]
  \item \textbf{Master (rank 0)}: genera lotes de trabajo y agrega resultados.
  \item \textbf{Workers (rank >= 1)}: reciben lotes de frames por MPI, ejecutan preprocesado CUDA e inferencia.
  \item \textbf{Salida}: escritura de CSV, JSON y video anotado ordenado por frame.
\end{enumerate}

\subsection{Planificaci\'on MPI}
Modos soportados:
\begin{itemize}[leftmargin=*]
  \item \texttt{dynamic}
  \item \texttt{static-contiguous}
  \item \texttt{static-round-robin}
\end{itemize}

El modo din\'amico mejora adaptaci\'on a variabilidad por frame. Los modos est\'aticos permiten comparar equilibrio de carga frente a una distribuci\'on fija.

\section{Dise\~no de m\'odulos}
\subsection{Entrada y configuraci\'on}
\begin{itemize}[leftmargin=*]
  \item \path{src/main.cpp}: parsing de argumentos, validaci\'on y bootstrap de entorno para bibliotecas CUDA/cuDNN.
\end{itemize}

\subsection{Orquestaci\'on distribuida}
\begin{itemize}[leftmargin=*]
  \item \path{src/mpi_runtime.cpp}: distribuci\'on de lotes, env\'io de cabeceras MPI, recepci\'on de resultados y merge ordenado.
\end{itemize}

\subsection{Lectura de video}
\begin{itemize}[leftmargin=*]
  \item \path{src/video_io.cpp}: lectura por lotes desde frame inicial arbitrario en el master para alimentar a los workers.
\end{itemize}

\subsection{Preprocesado e inferencia (Polimorfismo)}
Para dar soporte a la heterogeneidad del hardware, la l\'ogica de detecci\'on est\'a abstra\'ida mediante la interfaz \texttt{IPeopleDetector}.
\begin{itemize}[leftmargin=*]
  \item \textbf{Backend CUDA}:
  \begin{itemize}[leftmargin=*]
    \item \path{src/gpu_preprocess.cu}: kernels personalizados en CUDA C++ para letterboxing y filtros Sobel ultrarr\'apidos.
    \item \path{src/yolo_detector_cuda.cpp}: inferencia de m\'aximo rendimiento usando la API C++ de ONNX Runtime con el proveedor de ejecuci\'on CUDA.
  \end{itemize}
  \item \textbf{Backend OpenCL}:
  \begin{itemize}[leftmargin=*]
    \item \path{src/gpu_preprocess_opencl.cpp}: preprocesado hardware-agnostic basado en la abstracci\'on \texttt{cv::UMat} de OpenCV, manteniendo la memoria en la VRAM de la gr\'afica.
    \item \path{src/yolo_detector_opencl.cpp}: inferencia compatible con GPUs integradas Intel/AMD apoy\'andose en el m\'odulo \texttt{dnn} nativo de OpenCV (\texttt{DNN\_TARGET\_OPENCL}).
  \end{itemize}
\end{itemize}

\subsection{Postproceso y salida}
\begin{itemize}[leftmargin=*]
  \item \path{src/tracker.cpp}: suavizado de cajas para anotaci\'on visual.
  \item \path{src/output_writer.cpp}: CSV, JSON y video final anotado.
\end{itemize}

\section{Soporte H\'ibrido y Optimizaci\'on OpenCL}
Uno de los principales hitos t\'ecnicos de este analizador es su capacidad para orquestar recursos de c\'omputo dispares.

\subsection{Despacho Din\'amico Inteligente}
En \path{src/main.cpp} y \path{src/mpi_runtime.cpp}, el sistema inspecciona autom\'aticamente la topolog\'ia del nodo. Si el usuario utiliza el modo \texttt{--gpu-backend auto}, el programa:
\begin{enumerate}
    \item Cuenta las GPUs de NVIDIA detectables (ej. 1 GPU RTX).
    \item Asigna prioritariamente los primeros \emph{workers} de MPI a estas GPUs v\'ia CUDA.
    \item Cualquier \emph{worker} extra instanciado (ej. \emph{Rank 2}) detecta que las GPUs dedicadas est\'an saturadas y hace un \emph{fallback} elegante hacia el backend de OpenCL, adue\~n\'andose de la GPU integrada.
\end{enumerate}

\subsection{Evasi\'on de Cuellos de Botella de Memoria (UMat)}
Durante la implementaci\'on inicial de OpenCL, se observ\'o un estancamiento en la GPU integrada (con usos en torno al 12\%). El an\'alisis revel\'o que la funci\'on \texttt{blobFromImage} descargaba s\'incronamente los tensores desde la VRAM a la RAM de la CPU (como un \texttt{cv::Mat}), para que acto seguido la red neuronal volviera a subirlos a la GPU. 

La soluci\'on consisti\'o en forzar la salida de dicha funci\'on hacia un objeto \texttt{cv::UMat} unificado. Al mantener el tensor confinado en la memoria de v\'ideo de principio a fin, el tiempo de inferencia de la integrada se redujo a una tercera parte, estabilizando el flujo as\'incrono.

\section{Planificaci\'on en Entornos Asim\'etricos}
Para sacar verdadero partido a una GPU dedicada (NVIDIA RTX) junto con una GPU integrada simult\'aneamente, el planificador \texttt{--scheduler dynamic} result\'o ser matem\'aticamente superior.
Mientras la NVIDIA procesa un lote en 20 ms, la integrada requiere unos 200 ms. Si se utilizara un planificador \texttt{static-round-robin} estricto, la tarjeta r\'apida terminar\'ia su mitad de lotes del v\'ideo de forma casi instant\'anea y esperar\'ia ociosa a que la lenta terminase. El modo din\'amico soluciona esto alimentando vorazmente a la GPU r\'apida y dejando que la integrada procese tranquilamente de fondo, logrando un ratio natural de procesamiento de 10 a 1 sin bloqueos por sincronizaci\'on ineficiente.

\section{Formato de datos y trazabilidad}
\subsection{CSV por frame}
Columnas principales:
\begin{itemize}[leftmargin=*]
  \item \texttt{frame\_index, worker\_rank, timestamp\_ms}
  \item \texttt{people\_count}
  \item \texttt{read\_ms, preprocess\_ms, detection\_ms, processing\_ms}
  \item \texttt{people\_count} se limita al n\'umero de cajas transmitidas (tope \texttt{kMaxDetectionsPerFrame}).
\end{itemize}

\subsection{Resumen JSON}
Incluye metadatos de video, m\'etricas agregadas, balance de carga por worker y tiempos globales de proceso.

\section{Repaso t\'ecnico del estado actual}
\subsection{Fortalezas}
\begin{itemize}[leftmargin=*]
  \item Arquitectura MPI coherente con reparto por lotes serializados desde el master, sin relectura del video en cada worker.
  \item Separaci\'on clara de tiempos por etapa, \`util para an\'alisis de cuellos de botella.
  \item Salida reproducible y auditable en CSV + JSON + video.
  \item Backend de inferencia alineado con GPU para el caso principal.
\end{itemize}

\subsection{Aspectos a vigilar}
\begin{itemize}[leftmargin=*]
  \item La contenci\'on de una GPU compartida puede limitar escalado al aumentar workers.
  \item El tracker no debe interpretarse como m\'odulo de tracking de identidades persistentes.
  \item El acumulado visual representa una estimaci\'on de personas \`unicas basada en IDs de tracking.
\end{itemize}

\section{Mejora implementada: contador actual y acumulado}
\subsection{Motivaci\'on}
Se solicit\'o mostrar en el video procesado:
\begin{enumerate}[leftmargin=*]
  \item cu\'antas personas hay en el frame actual,
  \item y cu\'al es el conteo acumulado hasta ese instante.
\end{enumerate}

\subsection{Implementaci\'on}
En la etapa de anotaci\'on de \texttt{output\_writer}:
\begin{itemize}[leftmargin=*]
  \item se reconstruyen tracks visuales sobre las detecciones ordenadas por frame,
  \item se obtiene \texttt{current\_people} como el n\'umero de tracks visibles,
  \item se mantiene un conjunto de IDs ya vistos para estimar \texttt{cumulative\_people},
  \item se dibuja un overlay en la esquina superior izquierda con ambos valores,
  \item se mantiene el dibujo de cajas de detecci\'on como antes.
\end{itemize}

Se a\~nadieron adem\'as los campos \texttt{accumulated\_people} y \texttt{accumulated\_detections} al \texttt{summary.json} para distinguir entre estimaci\'on de personas \`unicas y suma temporal de detecciones.

\section{Entorno de desarrollo y compilaci\'on}
\subsection{Dependencias}
\begin{itemize}[leftmargin=*]
  \item CMake
  \item OpenCV 4 (core, imgproc, videoio, dnn)
  \item OpenMPI
  \item OpenCL ICD (\texttt{intel-opencl-icd}, \texttt{mesa-opencl-icd})
  \item CUDA Toolkit (Opcional, manejado por CMake)
  \item ONNX Runtime GPU en entorno Python local (Opcional)
\end{itemize}

\subsection{Compilaci\'on}
El sistema de construcci\'on en \texttt{CMakeLists.txt} ha sido modularizado. Si el sistema carece de aceleradores NVIDIA, la opci\'on \texttt{RESCUE\_ENABLE\_CUDA=OFF} permite compilar la aplicaci\'on estrictamente para procesadores e integradas OpenCL, sin requerir el SDK de CUDA.
\begin{lstlisting}[language=bash]
cmake -S . -B build -DRESCUE_ENABLE_CUDA=ON
cmake --build build -j
\end{lstlisting}

\subsection{Ejecuci\'on recomendada (Ejemplo H\'ibrido)}
Para aprovechar simult\'aneamente una GPU NVIDIA y una GPU Integrada Intel, deben lanzarse 3 procesos (1 Master + 2 Workers) con el balanceador auto-configurado:
\begin{lstlisting}[language=bash]
LD_LIBRARY_PATH=/usr/lib/wsl/lib:${LD_LIBRARY_PATH:-} \
mpiexec -n 3 ./build/rescue_video_analyzer \
  --input dataset/videoset2.mp4 \
  --output-dir output \
  --batch-size 8 \
  --scheduler dynamic \
  --gpu-backend auto
\end{lstlisting}

\section{Resultados y lectura de m\'etricas}
\subsection{M\'etricas clave}
\begin{itemize}[leftmargin=*]
  \item tiempo medio de lectura por frame,
  \item tiempo medio de preprocesado,
  \item tiempo medio de detecci\'on,
  \item tiempo total por frame,
  \item personas promedio por frame y m\'aximo observado,
  \item acumulado total del video.
\end{itemize}

\subsection{Interpretaci\'on}
\begin{itemize}[leftmargin=*]
  \item Si \texttt{detection\_ms} domina, la inferencia es cuello de botella.
  \item Si \texttt{read\_ms} sube con muchos workers, puede haber contenci\'on de E/S.
  \item Si el balance de carga muestra \texttt{max\_to\_mean\_ratio} alto, conviene probar planificador din\'amico o tama\~no de lote distinto.
\end{itemize}

\section{Validaci\'on funcional recomendada}
\begin{enumerate}[leftmargin=*]
  \item Ejecutar un video corto y verificar que se genera \texttt{annotated\_people.mp4} o \texttt{.avi}.
  \item Confirmar que el overlay muestra ambos contadores y crece el acumulado.
  \item Comparar \texttt{accumulated\_people} en JSON con el comportamiento visual del tracker y usar \texttt{accumulated\_detections} si quieres la suma temporal de \texttt{people\_count}.
  \item Revisar logs de workers para backend y latencias por lote.
\end{enumerate}

\section{Gu\'ia de compilaci\'on de esta memoria}
Si hay distribuci\'on LaTeX instalada:
\begin{lstlisting}[language=bash]
pdflatex memoria_proyecto.tex
pdflatex memoria_proyecto.tex
\end{lstlisting}

Se recomienda doble compilaci\'on para actualizar tabla de contenidos y referencias.

\section{Trabajo futuro}
\begin{itemize}[leftmargin=*]
  \item Conteo acumulado por identidades \`unicas (tracking robusto por ID).
  \item Evaluaci\'on sistem\'atica de sensibilidad a umbrales de confianza y NMS.
  \item Exportaci\'on de reportes autom\'aticos (tablas y gr\'aficas) desde los CSV.
  \item Soporte multi-GPU real para escalado horizontal de workers.
\end{itemize}

\section{Conclusiones}
El proyecto presenta una base t\'ecnica s\'olida para an\'alisis de personas en video con procesamiento distribuido y aceleraci\'on GPU. La mejora implementada del overlay de conteo actual y acumulado aumenta el valor explicativo del video final y la trazabilidad del resultado. La memoria deja documentados alcance, arquitectura, flujo, m\'etricas y l\'ineas de evoluci\'on para siguientes iteraciones.

\end{document}
